\section{Auditoría de fuentes y mapa de lectura: qué responde realmente esta clase}

Esta clase se grabó el 24 de enero de 2023. Jason Wei conecta, mediante una cadena didáctica continua, dos temas de investigación que entonces acababan de formarse. Primero: ¿por qué la language-model loss puede disminuir de forma suave con la escala, mientras que ciertas capacidades en tareas parecen surgir de golpe, como si cruzaran un umbral? Segundo: ¿por qué la chain-of-thought (CoT, cadena de pensamiento) no funciona en todos los modelos, sino que produce una ganancia significativa solo cuando el modelo ya tiene suficiente capacidad básica? En esta reescritura, las 37 páginas oficiales de Google Slides son el eje visual. Las páginas 1--36 contienen material didáctico, y la página 37 es solo de agradecimientos. La comparación en vivo en OpenAI Playground entre las respuestas sin CoT y con CoT se recuperó extrayendo fotogramas exactos de la grabación oficial.

\lecturefigure{slide-01.jpg}{Tema de la clase: Scaling unlocks emergent abilities in language models}{Página 1 de las diapositivas oficiales}

La palabra ``unlock'' del título es importante. No afirma que el número de parámetros cree por sí mismo un módulo misterioso. Lo que dice es que, cuando el modelo, los datos, el cómputo de entrenamiento y la forma de uso entran juntos en cierta región, empiezan a ser medibles capacidades que antes no se veían en las métricas de la tarea. Conviene leer esta clase como un curso sobre «cómo interpretar curvas» y no tomar ``emergence'' como un término de mercadotecnia.

\lecturefigure{slide-02.jpg}{Ruta de la clase: capacidades emergentes, U-shaped scaling, CoT, BIG-Bench, razonamiento entre idiomas y self-consistency}{Página 2 de las diapositivas oficiales}

Las diapositivas tratan primero la relación entre la escala del preentrenamiento y la capacidad en tareas, y después pasan a las intervenciones de razonamiento en prompting-time. Las dos partes no son reseñas paralelas. La CoT es precisamente un caso clave: muestra que un método puede ser inútil, o incluso perjudicial, en modelos pequeños y, en cambio, abrir nuevas regiones de tareas en modelos grandes. La self-consistency muestra además que aumentar el presupuesto de muestreo durante la inferencia también tiene un umbral de escala.

\begin{importantbox}{La pregunta central de esta clase}
Dada una curva de escala del modelo contra desempeño, ¿cómo distinguimos entre una mejora suave, un umbral de medición, un cambio real de mecanismo y la ganancia que aportan las técnicas de prompting? Más aún, si los datos, el postentrenamiento, la metric y el inference compute pueden desplazar la curva, ¿qué magnitud debe ir en el eje horizontal?
\end{importantbox}

\teachervoice{Al inicio, el ponente subraya que una de las razones por las que este trabajo es importante es que investigadores de Google, DeepMind y Stanford construyeron un marco de observación común. Ante todo, es un lenguaje que permite el diálogo entre distintos modelos, tareas y equipos, y no una ley natural sobre un umbral fijo de parámetros.}

\subsection{Límites de la evidencia y principios de la reescritura}

Estas notas reconstruyen estrictamente la argumentación anterior a la fecha de la clase. La página oficial del curso recomienda tres artículos: Emergent Abilities, Chain-of-Thought Prompting y Scaling Instruction-Finetuned Language Models. Las diapositivas también usan directamente resultados de primera mano como scaling laws, BIG-Bench, BIG-Bench Hard, multilingual CoT y self-consistency. El debate que surgió más adelante en 2023 sobre si «la emergencia es un metric artifact» tiene valor para la investigación, pero no forma parte del contenido de la clase del 24 de enero, así que no se presentará retroactivamente como conclusión del ponente en ese momento.

\begin{warningbox}{El material histórico no debe reescribirse con información posterior}
En esta clase, ``emergent'' es un término operativo: los modelos pequeños quedan cerca del azar en la metric elegida y los modelos grandes lo superan con claridad. No equivale automáticamente a que haya aparecido dentro del modelo un nuevo circuito discreto, ni demuestra que exista una relación causal única entre el umbral y los parámetros del modelo. Estas notas conservan la evidencia presentada en clase y, al mismo tiempo, señalan con claridad las limitaciones que introducen la metric, los datos y la familia de modelos.
\end{warningbox}

\subsection{Resumen de la sección}

Al leer las gráficas que siguen, hay que tener presentes tres preguntas: qué es exactamente el eje horizontal; si el eje vertical tiene una baseline aleatoria o humana; y si el punto de inflexión de la curva proviene de la capacidad del modelo, la calidad de los datos, el postentrenamiento, el prompt o de la acción conjunta de estos factores con la metric. Todos los juicios posteriores sobre «emergencia» deben volver a estas tres preguntas.
